% VibeCFD — Journal Article Template
% Compatible with most Elsevier, Springer, and MDPI journals
\documentclass[preprint,12pt]{elsarticle}

\usepackage{amsmath,amssymb,amsfonts}
\usepackage{graphicx}
\usepackage{subcaption}
\usepackage{booktabs}
\usepackage{hyperref}
\usepackage{siunitx}
\usepackage{cleveref}
\usepackage{pgfplots}
\pgfplotsset{compat=1.18}

% SI units setup
\sisetup{
  per-mode=symbol,
  inter-unit-product=\ensuremath{{}\cdot{}},
}

\journal{Journal of Computational Physics}

\begin{document}

\begin{frontmatter}

\title{%% TITLE: Auto-generated by VibeCFD %%}

\author[inst1]{%% AUTHOR NAME %%}
\ead{%% EMAIL %%}
\affiliation[inst1]{organization={%% INSTITUTION %%}, city={%% CITY %%}, country={%% COUNTRY %%}}

\begin{abstract}
%% ABSTRACT: 150-250 words. State the problem, method, key results, and conclusions. %%
\end{abstract}

\begin{keyword}
computational fluid dynamics \sep OpenFOAM \sep %% KEYWORDS %%
\end{keyword}

\end{frontmatter}

%% ─── Introduction ─────────────────────────────────────────────
\section{Introduction}
\label{sec:introduction}

%% Background, literature review, research gap, and objectives. %%

%% ─── Mathematical Formulation ─────────────────────────────────
\section{Mathematical Formulation}
\label{sec:formulation}

The governing equations are the incompressible Navier-Stokes equations:
\begin{equation}
  \nabla \cdot \mathbf{u} = 0
  \label{eq:continuity}
\end{equation}
\begin{equation}
  \frac{\partial \mathbf{u}}{\partial t} + (\mathbf{u} \cdot \nabla)\mathbf{u}
  = -\frac{1}{\rho}\nabla p + \nu \nabla^2 \mathbf{u}
  \label{eq:momentum}
\end{equation}
where $\mathbf{u}$ is the velocity vector, $p$ is pressure, $\rho$ is density,
and $\nu$ is kinematic viscosity.

%% Add turbulence model equations if applicable. %%

%% ─── Numerical Method ─────────────────────────────────────────
\section{Numerical Method}
\label{sec:method}

\subsection{Solver and Discretization}

%% AUTO-FILLED: Solver name, spatial/temporal discretization schemes. %%

\subsection{Computational Domain and Mesh}

%% AUTO-FILLED: Domain dimensions, mesh type, cell count, boundary conditions. %%

\begin{table}[htbp]
  \centering
  \caption{Mesh parameters and quality metrics.}
  \label{tab:mesh}
  \begin{tabular}{lrrr}
    \toprule
    Parameter & Coarse & Medium & Fine \\
    \midrule
    Cells              & %% N1 %% & %% N2 %% & %% N3 %% \\
    Max non-orthogonality (\si{\degree}) & %% %% & %% %% & %% %% \\
    Max skewness       & %% %% & %% %% & %% %% \\
    Mean $y^+$         & %% %% & %% %% & %% %% \\
    \bottomrule
  \end{tabular}
\end{table}

\subsection{Boundary Conditions}

\begin{table}[htbp]
  \centering
  \caption{Boundary conditions.}
  \label{tab:bc}
  \begin{tabular}{llll}
    \toprule
    Boundary & Type & $\mathbf{u}$ & $p$ \\
    \midrule
    Inlet    & Dirichlet & %% %% & zeroGradient \\
    Outlet   & Neumann   & zeroGradient & %% %% \\
    Walls    & No-slip   & $(0,0,0)$ & zeroGradient \\
    \bottomrule
  \end{tabular}
\end{table}

%% ─── Results ──────────────────────────────────────────────────
\section{Results and Discussion}
\label{sec:results}

\subsection{Mesh Independence Study}

%% AUTO-FILLED: GCI table, Richardson extrapolation. %%

\subsection{Validation}

%% AUTO-FILLED: Comparison with benchmark data. %%

\begin{figure}[htbp]
  \centering
  %% \includegraphics[width=\columnwidth]{figures/validation.pdf} %%
  \caption{Comparison of present results with benchmark data.}
  \label{fig:validation}
\end{figure}

\subsection{Flow Analysis}

%% AUTO-FILLED: Velocity contours, streamlines, pressure distribution. %%

%% ─── Conclusions ──────────────────────────────────────────────
\section{Conclusions}
\label{sec:conclusions}

%% Summary of findings, contributions, and future work. %%

%% ─── References ───────────────────────────────────────────────
\bibliographystyle{elsarticle-num}
\bibliography{references}

\end{document}
